We have introduced partisan similarity districts (PSDs): algorithmic maps
generated by weighting census-tract adjacency edges by partisan vote
similarity rather than shared boundary length.  PSDs produce within-district
partisan homogeneity substantially exceeding enacted gerrymanders---mean
within-district $\sigma_D$ falls from 14.8 to 7.8 pp in NC and from 16.2
to 9.2 pp in TX at $\alpha = 50$---while distributing safe seats
proportionally between parties (EG $\approx 0$).\singrun

\paragraph{Key finding.}
Enacted partisan gerrymanders operate at an implicit partisan-similarity
level equivalent to approximately $\alpha \approx 15$--$20$ in our
parameterisation, but do so in a directionally asymmetric way: they generate
excess safe seats for the mapmaking party beyond what symmetric geographic
clustering would produce.  The NC enacted plan has 3 more safe Republican
seats than a neutral PSD at $\alpha = 20$ and 3 fewer safe Democratic seats,
with an EG difference of $0.15$ over the PSD baseline.\singrun

\paragraph{Legal utility.}
PSDs provide a quantitative answer to the \emph{Rucho}-era question: how
much partisan manipulation is embedded in an enacted plan, beyond what
geography requires?  The PSD safe-seat count is the geographic baseline;
the enacted-plan count is the observed outcome; the difference is the
manipulation premium.  Courts, commissions, and legislatures evaluating
redistricting plans can use this comparison without evaluating the
normatively contested question of whether safe seats are desirable.

\paragraph{Future work.}
Three extensions are natural.  First, ensemble analysis (multiple seeds and
METIS calls) would characterise the distribution of PSD maps and establish
confidence intervals on safe-seat counts.  Second, combining partisan-
similarity and VRA-aligned edge weights would test whether safe
minority-opportunity seats can be generated simultaneously with safe partisan
seats.  Third, extending PSD analysis to state legislative maps would cover
the 94\% of U.S. legislative seats that are not congressional, where partisan
gerrymandering effects are arguably even larger.
